\chapter{Glossaire}

\begin{longtable}[c]{| m{3.6cm} | m{11.4cm} |}
\caption{Glossaire des termes techniques}\\
\hline
\textbf{Terme} & \textbf{Définition} \\
\hline
\endfirsthead
\hline
\textbf{Terme} & \textbf{Définition} \\
\hline
\endhead
\hline
Abstention & Comportement par lequel le système refuse explicitement de répondre lorsque l'information demandée ne figure pas dans le corpus. Considéré ici comme un succès, non comme un échec. \\
\hline
Ancrage (\textit{grounding}) & Fait de contraindre une réponse générée à s'appuyer uniquement sur des documents fournis, et non sur les connaissances internes du modèle. \\
\hline
Base vectorielle & Base de données spécialisée dans le stockage de vecteurs et la recherche des plus proches voisins. \\
\hline
\textit{Chunk} & Unité de découpage du corpus. Ici, un chunk correspond exactement à un article de loi. \\
\hline
Cross-encodeur & Modèle qui évalue conjointement une paire (question, passage) pour en estimer la pertinence. Plus précis mais plus coûteux qu'un calcul de similarité entre vecteurs. \\
\hline
Embedding & Représentation vectorielle d'un texte, construite de sorte que la proximité géométrique traduise la proximité de sens. \\
\hline
Hallucination & Production par un modèle de langage d'un contenu plausible mais faux --- typiquement, en contexte juridique, un numéro d'article inexistant. \\
\hline
MRR & \textit{Mean Reciprocal Rank} : moyenne de l'inverse du rang du premier résultat pertinent. Mesure la capacité à placer le bon résultat en tête. \\
\hline
Quantification & Réduction de la précision numérique des poids d'un modèle, afin de diminuer son empreinte mémoire et de permettre son exécution sur du matériel ordinaire. \\
\hline
RAG & \textit{Retrieval-Augmented Generation} : architecture où la génération est précédée d'une recherche documentaire, dont les résultats sont fournis au modèle. \\
\hline
Recall@k & Proportion de questions pour lesquelles le document attendu figure parmi les $k$ premiers résultats. \\
\hline
Similarité cosinus & Mesure de proximité entre deux vecteurs, égale au cosinus de l'angle qu'ils forment. Bornée entre $-1$ et $1$. \\
\hline
Souveraineté (des données) & Principe selon lequel une donnée produite sur un territoire doit pouvoir y être traitée, sans dépendance à une infrastructure étrangère. \\
\hline
Température & Paramètre de génération contrôlant l'aléatoire du modèle. Une valeur basse produit des réponses plus stables et reproductibles. \\
\hline
\end{longtable}
